\section{From subgroup congruences to procyclic length}\label{sec:profinite}

We now isolate the group-theoretic consequence of the character calculations.
All homomorphisms between profinite groups are understood to be continuous, and a finite
section means a finite quotient of a closed subgroup.  For a finite
$p$-group $F$, write $d(F)$ for its minimum number of generators.

\begin{theorem}\label{thm:procyclic}
Let $\Gamma$ be a finitely generated pro-$p$ group, and let $l\geq0$.
Suppose that $d(F)\leq l$ for every finite section $F$ of $\Gamma$, that
\begin{equation}\label{eq:profinite-hom}
 \frac{|\Hom_{\cts}(\Gamma,F)|}{|F|}\in\Z_{(p)}
 \qquad\text{for every finite group }F,
\end{equation}
and that, for every $U\leq_o\Gamma$, every finite continuous $U$-set $X$,
and every $h\geq0$,
\begin{equation}\label{eq:profinite-congruences}
 \sum_{\substack{V\leq_o U\\{}[U:V]=p^h}}
       \bigl(|X^V|-|X|\bigr)
       \in p^{h+1}\Z_{(p)}.
\end{equation}
Then $\Gamma$ has a series of closed subgroups
\[
 1=K_t\trianglelefteq K_{t-1}\trianglelefteq\cdots
   \trianglelefteq K_0=\Gamma,
 \qquad t\leq l,
\]
whose nontrivial factors are isomorphic to $\Z_p$.
In particular, every finite quotient $P$ of $\Gamma$ satisfies
$\ell(P,1)\leq l$.
\end{theorem}

The congruence in \eqref{eq:profinite-congruences} is summed over actual
open subgroups.  These sums are finite: an open subgroup of a given
index is a point stabilizer in a permutation action of that degree,
and a finitely generated profinite group has only finitely many such
actions.  We say that a finitely generated pro-$p$ group
\emph{satisfies the subgroup congruences} when this formula holds
for all its open subgroups.  This property is inherited by open
subgroups. For nontrivial $\Gamma$, homomorphism-count integrality
supplies the first $\Z_p$ quotient. After each quotient is extracted,
shrinking its action replaces a semidirect product by a direct product
while preserving the subgroup congruences and section bound.
Subgroup counting then gives another $\Z_p$ quotient whenever the
remaining kernel is nontrivial.

\subsection{Finite quotients and shrinking actions}

We first record the elementary profinite facts needed to make that
limit precise.

\begin{lemma}\label{lem:profinite-elementary}
Let $E$ be a pro-$p$ group.
\begin{enumerate}[label=\textup{(\roman*)}]
\item Every finite section of $E$ is a section of a finite quotient of
$E$.  If all finite sections of $E$ have generator number at most $l$,
then every closed subgroup of $E$ is topologically generated by at most
$l$ elements.
\item A procyclic pro-$p$ group is either trivial, finite cyclic, or
isomorphic to $\Z_p$.  Every continuous epimorphism $K\twoheadrightarrow
\Z_p$ splits.
\item A finitely generated pro-$p$ group with infinite abelianization
has a continuous epimorphism onto $\Z_p$.
\end{enumerate}
\end{lemma}

\begin{proof}
For a closed subgroup $H\leq E$ and $N\trianglelefteq_o H$, choose
$M\trianglelefteq_o E$ with $H\cap M\subseteq N$.  Then $H/N$ is a
quotient of the subgroup $H/(H\cap M)$ of $E/M$.  Under the asserted
generator bound, for each $N\trianglelefteq_o H$ the set of tuples in
$H^l$ generating $H/N$ is nonempty and closed.  These sets have the
finite intersection property, by replacing finitely many $N$ by their
intersection.  Compactness supplies a tuple generating every finite
quotient of $H$, hence topologically generating $H$.

A topological generator gives a continuous epimorphism
$\Z_p\twoheadrightarrow Q$ onto any procyclic pro-$p$ group $Q$.
Every closed subgroup of $\Z_p$ is a $\Z_p$-submodule and is either
$0$ or $p^a\Z_p$, proving the classification.  If
$K\twoheadrightarrow\Z_p$, choose $x\in K$ mapping to $1$.
The closed subgroup generated by $x$ maps isomorphically to $\Z_p$:
a nonzero closed kernel in a procyclic group would make the quotient
finite.  This supplies a continuous section.

Finally, the abelianization of a finitely generated pro-$p$ group is a
finitely generated $\Z_p$-module.  The elementary structure theorem over
the discrete valuation ring $\Z_p$ writes it as a finite group times
$\Z_p^b$.  If it is infinite, then $b>0$.
\end{proof}

\begin{lemma}[Shrinking an action]\label{lem:shrinking-action}
Let $K$ be a pro-$p$ group and let $d,l\geq0$.
Suppose $K\times\Z_p^d$ satisfies the subgroup congruences and all its
finite sections have generator number at most $l$.  If
$K\twoheadrightarrow\Z_p$ has kernel $N$, then
$N\times\Z_p^{d+1}$ has both properties.
\end{lemma}

\begin{proof}
By \Cref{lem:profinite-elementary}, write
$K=N\rtimes_\alpha\Z_p$, with $N$ finitely generated.  For $s\geq0$,
put
\[
 E_s=(N\rtimes_{\alpha_s}\Z_p)\times\Z_p^d,
 \qquad \alpha_s(z)=\alpha(p^s z),
 \qquad E_\infty=N\times\Z_p^{d+1}.
\]
The group $E_s$ is isomorphic to the open subgroup
$(N\rtimes p^s\Z_p)\times\Z_p^d$ of $K\times\Z_p^d$, and hence has
both stated properties.

We construct a common finite quotient which captures every permutation
action of a prescribed bounded degree.  Fix $M\geq1$, and let
\[
 O_M=\bigcap_{\rho\in\Hom_{\cts}(N,\Sigma_M)}\ker\rho.
\]
There are finitely many such homomorphisms, because $N$ is finitely
generated.  Thus $O_M$ is open and characteristic in $N$.  Choose
$a\geq0$ with $p^a\geq M$.  Every $p$-element of $\Sigma_M$ has order
at most $p^a$, as its order is the largest length of one of its
$p$-power cycles.  Therefore every homomorphism from $\Z_p$ to
$\Sigma_M$ kills $p^a\Z_p$.

The action $\alpha$ on the finite group $N/O_M$ is continuous, so
$\alpha(p^s\Z_p)$ acts trivially on it for all sufficiently large $s$.
For these $s$, both $E_s$ and $E_\infty$ have the quotient
\begin{equation}\label{eq:shared-permutation-quotient}
 Q_M=(N/O_M)\times(\Z/p^a)^{d+1}.
\end{equation}
Every homomorphism from either group to $\Sigma_M$ factors through this
quotient: its restriction to $N$ kills $O_M$, and its restrictions to
the $d+1$ cyclic coordinates kill $p^a\Z_p$.
The same assertion holds for every smaller permutation degree, by
adding fixed points.

Now fix $U\leq_o E_\infty$, a finite continuous $U$-set $X$, and
$h\geq0$.  If $X$ is empty, the required congruence is immediate.
Otherwise set $m=[E_\infty:U]$ and take
\[
 M\geq\max\{m,\,m|X|,\,mp^h\}.
\]
The coset action on $E_\infty/U$ factors through $Q_M$, so $U$ is the
inverse image of a subgroup $\overline U\leq Q_M$.  The action on the
induced finite set $E_\infty\times_U X$ also factors through $Q_M$.
In particular, the quotient kernel acts trivially on its fiber $X$
over $U$, making $X$ a $\overline U$-set.  Let $U_s\leq E_s$ be the
inverse image of $\overline U$, and give the same set $X$ the resulting
$U_s$-action.

Every open subgroup $V\leq U$ of index $p^h$ has index $mp^h$ in
$E_\infty$.  Its coset action therefore factors through $Q_M$, so $V$
contains the quotient kernel.  Exactly the same argument applies to
every open subgroup of index $p^h$ in $U_s$.  Thus the two collections
of subgroups are both in bijection with
\[
 \{\overline V\leq\overline U:
                 [\overline U:\overline V]=p^h\}.
\]
The corresponding fixed-point counts on $X$ agree term by term.
The congruence for $U_s$ consequently proves the congruence for $U$.
This proves the required assertion for every open subgroup of
$E_\infty$ and every test on it.

For the generator bound, let $F$ be a finite quotient of $E_\infty$.
Its regular permutation action shows, using $M\geq|F|$, that the
quotient map factors through \eqref{eq:shared-permutation-quotient}.
Thus $F$ is also a quotient of $E_s$ for sufficiently large $s$.
Every finite section of $E_\infty$ is a section of such an $F$ by
\Cref{lem:profinite-elementary}; it therefore has generator number at
most $l$.
\end{proof}

\subsection{Detecting the remaining kernel}

The subgroup-counting formula below is classical; see
Tamanoi \cite[Lemma~5-5]{Tamanoi2001}. We include its proof and record
the \(p\)-adic limit needed for the kernel argument.

\begin{lemma}\label{lem:lattice-limit}
For $d\geq1$, let $a_j(d)$ be the number of open subgroups of index
$p^j$ in $\Z_p^d$.  Then
\begin{equation}\label{eq:lattice-generating-function}
 \begin{aligned}
 \sum_{j\geq0}a_j(d)t^j
       &=\prod_{u=0}^{d-1}(1-p^u t)^{-1},\\
 \lim_{j\to\infty}a_j(d)
       &=c_d:=\prod_{u=1}^{d-1}(1-p^u)^{-1}\ne0
 \quad\text{in }\Z_p.
 \end{aligned}
\end{equation}
\end{lemma}

\begin{proof}
For $d=1$ there is exactly one subgroup of each index.
For $d>1$, project a sublattice of
$\Z_p\times\Z_p^{d-1}$ onto the first coordinate.  If its image is
$p^a\Z_p$ and its intersection with the remaining coordinates has
index $p^b$, it is determined by that intersection and a lift of
$p^a$ modulo it.  There are $p^b$ choices of lift, and the total index
is $p^{a+b}$.  Consequently
\[
 a_j(d)=\sum_{a+b=j}p^b a_b(d-1),
\]
which gives the displayed generating function by induction.

Write
\[
 \prod_{u=1}^{d-1}(1-p^u t)^{-1}
        =\sum_{m\geq0}b_m t^m.
\]
Each $b_m$ is divisible by $p^m$, and the generating function gives
$a_j(d)=\sum_{m=0}^j b_m$.  These partial sums converge $p$-adically.
The convergent geometric series in the finite product identify their
limit with $c_d$.  Every factor $1-p^u$, $u\geq1$, is a $p$-adic unit,
so $c_d$ is nonzero.  The empty product for $d=1$ is $1$.
\end{proof}

\begin{lemma}\label{lem:kernel-abelianization}
Let $K$ be a nontrivial finitely generated pro-$p$ group and $d\geq1$.
If $K\times\Z_p^d$ satisfies the subgroup congruences, then $K$ has
infinite abelianization.
\end{lemma}

\begin{proof}
Let $\Phi(K)$ be the intersection of the kernels of homomorphisms
$K\to C_p$.  There are finitely many such homomorphisms, so
$F=K/\Phi(K)$ is a finite elementary abelian group.  It is nontrivial:
a nontrivial pro-$p$ group has a nontrivial finite $p$-group quotient,
and every nontrivial finite $p$-group has a quotient $C_p$.
We will use the following precise property of $F$:
\begin{equation}\label{eq:frattini-detector}
 J\leq K\text{ closed and }J\twoheadrightarrow F
       \quad\Longrightarrow\quad J=K.
\end{equation}
Indeed, if $J$ were proper, its image would be proper in some finite
quotient of $K$.  A maximal subgroup containing that image is normal
of index $p$, as in every finite $p$-group.  Its inverse image is the
kernel of a map $K\twoheadrightarrow C_p$ containing $J$, contradicting
surjectivity onto $F$.

The rational full-mark idempotent $e_F^F$ is a rational virtual $F$-set
with mark $1$ at $F$ and mark $0$ at every proper subgroup.  Since
$F\ne1$, its underlying virtual cardinality is zero.  Choose a positive
integer $D$ clearing its denominators and inflate $De_F^F$ to
$E=K\times\Z_p^d$ through $E\to K\to F$.
For a closed subgroup of $E$, its projection to $K$ is closed.
By \eqref{eq:frattini-detector}, the inflated virtual set has mark $D$
exactly when that projection is all of $K$, and mark zero otherwise.

Let $c_h$ be the number of open subgroups of index $p^h$ in $E$
projecting onto $K$.  The congruences extend by additivity to integral
virtual sets.  Applying them to $De_F^F$ at $U=E$ gives
\begin{equation}\label{eq:surjective-count-vanishes}
 Dc_h\in p^{h+1}\Z_{(p)},
 \qquad v_p(c_h)\geq h+1-v_p(D).
\end{equation}
In particular, $c_h\to0$ in $\Q_p$.

We compute this count in a second way.  Set $A=\Z_p^d$.
For an open subgroup $H\leq K\times A$ projecting onto $K$, put
$V=H\cap K$ and $W=\operatorname{pr}_A(H)$.
Then $V\trianglelefteq_o K$: a lift in $H$ of any element of $K$
normalizes $V$.  Also $W\leq_o A$, and $H/V$ is the graph of a unique
continuous surjection
\[
 f:W\twoheadrightarrow K/V.
\]
Conversely, every such triple $(V,W,f)$ gives one subgroup $H$.
Its index is $[K:V][A:W]$.  Every open subgroup $W$ of $A$ is
isomorphic to $A$, and hence
\begin{equation}\label{eq:graph-subgroup-count}
 c_h=
 \sum_{\substack{V\trianglelefteq_o K\\{}[K:V]=p^i,\ i\leq h}}
       a_{h-i}(d)\,|\Sur_{\cts}(\Z_p^d,K/V)|.
\end{equation}
There is no automorphism denominator: the actual graph map is uniquely
determined by the actual subgroup.

Suppose now that $K^{\mathrm{ab}}$ were finite.  Any $K/V$ appearing
with a nonzero summand in \eqref{eq:graph-subgroup-count} is abelian,
so $V$ contains the closed commutator subgroup.  There are therefore
only finitely many contributing $V$.  By \Cref{lem:lattice-limit},
\[
 \lim_{h\to\infty}c_h
    =c_d\sum_{V\trianglelefteq_o K}
                   |\Sur_{\cts}(\Z_p^d,K/V)|\ne0,
\]
where the sum has only finitely many nonzero terms and is a positive
integer: $V=K$ contributes $1$.  This contradicts
\eqref{eq:surjective-count-vanishes}, proving the lemma.
\end{proof}

\subsection{The series and the rational gluing estimate}

\begin{proof}[Proof of \Cref{thm:procyclic}]
If $\Gamma=1$, the empty series suffices.  Otherwise $\Gamma$ has
infinite abelianization.  For if $\Gamma^{\mathrm{ab}}$ were finite
of order $p^b$, then for every sufficiently large $r$,
\[
 |\Hom_{\cts}(\Gamma,C_{p^r})|=p^b.
\]
This follows by decomposing the finite abelianization into cyclic
factors, each of which has as many maps to $C_{p^r}$ as it has
elements once $p^r$ is divisible by its exponent. Equation
\eqref{eq:profinite-hom} would then assert
$p^{b-r}\in\Z_{(p)}$ for arbitrarily large $r$, a contradiction.
Thus \Cref{lem:profinite-elementary} supplies an epimorphism
$K_0=\Gamma\twoheadrightarrow\Z_p$ with kernel $K_1$.

By \Cref{lem:shrinking-action}, $K_1\times\Z_p$ satisfies the
subgroup congruences and has the same finite-section bound.  More
generally, after $d$ steps we obtain $K_d\times\Z_p^d$ with these
properties.  If $K_d\ne1$, it is finitely generated by
\Cref{lem:profinite-elementary}, and for $d\geq1$,
\Cref{lem:kernel-abelianization} supplies an epimorphism
$K_d\twoheadrightarrow\Z_p$.  Take its kernel $K_{d+1}$ and apply
\Cref{lem:shrinking-action} again.

If $K_d\ne1$, it has a quotient $C_p$, so
$K_d\times\Z_p^d$ has a quotient $(C_p)^{d+1}$.  The section bound
therefore forces $d+1\leq l$.  The construction must terminate with
$K_t=1$ for some $t\leq l$, and
$K_{i-1}/K_i\cong\Z_p$ at every step.  Taking the images of these
subgroups in a finite quotient $P$ gives a subnormal series with
cyclic factors.  Deleting repetitions proves $\ell(P,1)\leq l$.
\end{proof}

\begin{corollary}[Rational gluing estimate]\label{cor:rational-gluing}
For every finite $p$-group $P$ and every $l\geq0$,
\[
 \bigl(U_P(L_{T(l)}S)\bigr)_{\Q}\ne0
       \quad\Longrightarrow\quad \ell(P,1)\leq l.
\]
\end{corollary}

\begin{proof}
For $l\geq1$, \Cref{thm:pro-characters} constructs a pro-$p$ group
$\Gamma\twoheadrightarrow P$ with the required finite-section bound.
\Cref{lem:hom-integrality,thm:subgroup-congruence} give
\eqref{eq:profinite-hom} and \eqref{eq:profinite-congruences}.
The conclusion follows from \Cref{thm:procyclic}.

For $l=0$, let $P\ne1$.  Rational group cohomology, computed by
averaging, gives $(H\Q)^{BP}\simeq H\Q$.
The transfer class of $P/1$ acts on these cochains as multiplication
by $|P|$, which is invertible, whereas its full $P$-mark is zero.
The full-mark summand is consequently zero.  The rational
categorical-to-geometric fixed-point comparison of
\Cref{lem:occurrence-fixed-points} gives
$\bigl(U_P(H\Q)\bigr)_{\Q}=0$.
Cofree completion and geometric fixed points preserve rational
objects, since multiplication by every nonzero integer remains an
equivalence.  Hence $U_P(H\Q)=0$.
Thus nonvanishing at $l=0$ forces $P=1$, as required.
\end{proof}

The next elementary consequence of inverse limits will subtract a
free abelian rank in the higher-height argument.

\begin{lemma}\label{lem:product-series}
Let $P$ be a finite $p$-group, and let $k,l\geq0$.
If
\[
 \ell\bigl(P\times(C_{p^r})^k,1\bigr)\leq l
          \qquad\text{for every }r\geq1,
\]
then $k+\ell(P,1)\leq l$.
\end{lemma}

\begin{proof}
Write $G_r=P\times(C_{p^r})^k$.
At each level let $\mathcal S_r$ be the finite nonempty set of all
cyclic-quotient subnormal chains from $1$ to $G_r$ of length exactly
$l$, allowing repetitions.  Projection induces maps
$\mathcal S_{r+1}\to\mathcal S_r$.
There is a compatible choice of chains: in the compact product of
the finite sets $\mathcal S_r$, any finite collection of compatibility
conditions is satisfied by projecting a chain from the largest level
mentioned.  The finite intersection property applies even though the
maps between the sets of chains need not be surjective.

Denote the selected chain at level $r$ by $(H_{i,r})_{0\leq i\leq l}$.
Compatibility means that every map $H_{i,r+1}\to H_{i,r}$ is
surjective.  The inverse limits give closed subgroups
\[
 1=H_0\trianglelefteq H_1\trianglelefteq\cdots
   \trianglelefteq H_l=P\times\Z_p^k.
\]
Moreover,
\[
 H_i/H_{i-1}\cong
       \varprojlim_r H_{i,r}/H_{i-1,r}.
\]
To see surjectivity of this comparison, take a compatible tuple of
cosets.  Its nonempty finite sets of representatives at the finite
levels have compatible choices by compactness, giving a representative
in $H_i$.  The kernel is $H_{i-1}$.
Each displayed inverse limit is procyclic: the surjective transition
maps of its cyclic finite quotients send generators to generators,
so compatible generators exist by the same compactness argument,
and their resulting closed cyclic subgroup surjects to every level.

Let $A_i$ be the image of $H_i$ in $\Z_p^k$, and set
$J_i=H_i\cap P$.  These subgroups are closed, and there is an exact
sequence
\begin{equation}\label{eq:product-series-exact}
 1\longrightarrow J_i/J_{i-1}
  \longrightarrow H_i/H_{i-1}
  \longrightarrow A_i/A_{i-1}\longrightarrow1.
\end{equation}
For the kernel assertion, an element of $H_i$ projecting into
$A_{i-1}$ can be multiplied by an element of $H_{i-1}$ with the same
projection; the resulting element belongs to $J_i$, and
$J_i\cap H_{i-1}=J_{i-1}$.

Each $A_i$ is a $\Z_p$-submodule of $\Z_p^k$.
The procyclic quotient $A_i/A_{i-1}$ has rank at most one, and
the rank increments add to $k$.  Exactly $k$ steps therefore increase
rank.  At any such step, $H_i/H_{i-1}$ is an infinite procyclic
pro-$p$ group, hence $\Z_p$ by \Cref{lem:profinite-elementary}.
Its finite subgroup $J_i/J_{i-1}$ is trivial.  Thus
$J_i=J_{i-1}$ at each of these $k$ steps.

The $J_i$ form a subnormal chain from $1$ to $P$.
Their factors embed in the procyclic factors of the $H_i$ and,
being finite, are cyclic.  Deleting the $k$ repetitions just found
leaves a cyclic-quotient chain of length at most $l-k$, proving the
claim.
\end{proof}
